% =====================================================================
%  SECTION 4 -- the negative survey
%  Twenty routes to uniformity in (L,a): one direction, three layers.
% =====================================================================
\clearpage
\part{The negative survey}
\section{Twenty routes, one direction}\label{sec:survey}

This section is the audit of what was tried and did not work. It is written to a single
standard: every route is stated as the theorem it actually proves, then killed by the
hypothesis that is false, and the failure is attributed to a file. The organisation is
forced by the material. There is one target (\S\ref{sec:survey-target}), one table of
twenty routes plus two by-products (\S\ref{sec:survey-table}), and one obstruction with
three mutually independent layers (\S\ref{sec:survey-three}). What survives is two items,
one of them a bug (\S\ref{sec:survey-survives}).

\subsection{The target}\label{sec:survey-target}

\begin{claimbox}{cN}
\claimhead{The target statement}{.} The uniformity asked for throughout the campaign is the
single sentence
\[
\boxed{\;\inf_{L,t}\, w_1\!\left(L,\,2.5,\,t\right) \;>\; 0\;}
\]
where $w_1=f_1^2/M_0$ is the atom fraction of the K\"all\'en--Lehmann measure
($f_1^2=\nu(\{m_1^2\})$, $M_0=\int d\nu=C(0)=\mathrm{Var}(O)$) for the fixed operator
$O=\tfrac14\mathrm{tr}\,F^2$. This is the gap of \cite{jaffe} read on the atom route of
\S\ref{sec:reduction-equiv}; the reading is not a weakening of the Clay statement, and
\S\ref{sec:survey-three} states exactly how far it is from it.
\src{\code{LOOP-RUNBOOK.md}:18982--18983 (\S160.1); definitions \code{LOOP-RUNBOOK.md}:18950--18956.}
\end{claimbox}

Three bookkeeping facts make this a well-posed sentence, and each was a separate round of
work to establish.

\begin{itemize}
\item \textbf{The lattice spacing is not a free parameter.} At the physical coupling
  $\beta=2.5$ the spacing is pinned, $a=a(\beta)$; the family is a graph, not a rectangle.
  Leaving $a$ free loses $\beta=2.5$; pinning it degrades ``$a$-uniformity'' to a
  statement about one point. Both readings are recorded in the corpus, and the second is
  the one used here.
\item \textbf{The ratio is scale invariant.} The engineering dimension of $w_1$ vanishes
  identically, $\gamma_{w_1}=2\gamma_O-2\gamma_O\equiv 0$, because numerator and
  denominator carry the same anomalous dimension. This is why no renormalisation of the
  bare ratio can move the target; the bare ratio nevertheless drifts to zero, $w_1^{\rm
  bare}\to0$, because the contact term makes $M_0\sim a^{-8}$. That drift is an artefact
  of not fixing the renormalisation scheme, not information about the atom.
\item \textbf{The order of the infima does not matter.} On a product set the three
  infima $\inf_L\inf_t$, $\inf_t\inf_L$ and $\inf_{(L,t)}$ coincide (a two-line lemma).
  ``The order is the difficulty'' is therefore false; what is genuinely open is the
  domain constraint, whether $\lim$ sits inside or outside the $\inf$, and the missing
  $t$-axis.
\end{itemize}

\begin{warnbox}
\textbf{The operator is not decoration.} $w_1$ depends on $O$, so \emph{every} infimum
above is taken at fixed $O$. The witness is one line and exact: on a fixed $H$ take the
legitimate sources $O_\theta=1+\varepsilon(\cos\theta\,\psi_1+\sin\theta\,\psi_2)$; then
$w_1=\cos^2\theta\in[0,1]$ while $\mathrm{Spec}(H)$ does not move at all. Hence $w_1$ is a
property of the pair $(H,O)$, never of $H$ alone, and a statement of the form ``$w_1>0$''
has no content until $O$ is named. This is the same fact that makes the measured factor
$20.2$ between two operators on one Markov chain a feature rather than an error.
\src{\code{U11.md} \S\S1--2, \code{LOOP-RUNBOOK.md}:19241--19252 and :19283--19289 (\S165).}
\end{warnbox}

\subsection{A table of twenty routes, one direction}\label{sec:survey-table}

Twenty routes were attacked, in the sense that each was pushed until it either produced a
uniform lower bound on $w_1$ or produced the exact reason it cannot. All twenty fail, and
they fail in one direction: nine give only upper bounds, nine are circular, two have
explicit counterexamples. Two further items were produced by the same sweep and are kept
because they are usable.

{\footnotesize
\setlength{\tabcolsep}{4pt}
\renewcommand{\arraystretch}{1.12}
\begin{longtable}{@{}>{\raggedright\arraybackslash}p{0.5cm} >{\raggedright\arraybackslash}p{3.15cm} >{\raggedright\arraybackslash}p{4.9cm} >{\raggedright\arraybackslash}p{1.15cm} >{\raggedright\arraybackslash}p{5.1cm}@{}}
\caption{Twenty routes to an $(L,a)$-uniform floor on $w_1$, grouped by the direction of what they
produce, with two by-products. The group headings partition the outputs, not the methods
(\S\ref{sec:survey-table}, notes).}\label{tab:routes}\\
\toprule
\textbf{\#} & \textbf{Route} & \textbf{What it gives} & \textbf{Grade} & \textbf{Cause of death}\\
\midrule
\endfirsthead
\toprule
\textbf{\#} & \textbf{Route} & \textbf{What it gives} & \textbf{Grade} & \textbf{Cause of death}\\
\midrule
\endhead
\bottomrule
\endlastfoot
\multicolumn{5}{@{}l}{\emph{Only upper bounds (9)}}\\
1 & compactness & Arzel\`a--Ascoli transports, does not produce: the output is the identity $\inf_K w_1=w_1(x_0)$ & \LN & $K$ is not compact ($a=0$ and $L=\infty$ have no member); relative compactness does not exclude $\ell=0$\\
2 & reflection positivity & $\widehat w_1(T)\downarrow\ \ge w_1$ \cite{os1973}, i.e.\ an upper bound; the direction is reversed & \LN & RP fixes the sign of the monotone quantity, only constrains $w_1$ from above; $2/(L-1)\searrow0$ and $w_1\equiv0$ both satisfy it\\
3 & Berezin--Lieb & $\sum_j\lambda_j\le\sum_j\varepsilon_j$ is a sum bound: $\lambda_1\le(1/N)\sum\lambda$ & \LN & a trace inequality bounds the bottom from above; its only lower-shaped corollary needs $u\ge0$, and on YM the hypothesis \emph{is} the target\\
4 & probabilistic reading \cite{chatterjee} & $w_1=P(S=m_1^2)$ and $|EX|^2\le E|X|^2$ is an upper bound & \LN & the modulus square is not the expectation of an indicator; $|\langle\mathbb{G}|\phi\rangle|^2$ is a difference of two positive currents, sign indefinite\\
5 & Krein--Rutman residual $D_\infty$ & information content a factor $2.0975$ (window $[0.690,1.448]$) & \LN & fails the free-field acid test: at $\beta=0$, $D_\infty=1$ while $w_1=0$\\
6 & concavity & the chord lemma is true: $f$ concave, $f\ge0$, left end $r_0$ $\Rightarrow$ $f(2.5)\ge r_0(\beta_2-2.5)/(\beta_2-\beta_1)$ & \LN & $w_1$ is not concave: the Lorentzian witness $v^2/(4v^2+D^2)$ changes the sign of $f''$ at $2v/\sqrt3$; ray concavity would force $w_1(\infty)\ge r_0$ against the free end\\
7 & RG modulus $\delta_A\cdot a$ & a pure number, $\delta_A\cdot a=0.5686159$ & \LN & a $0.24\%$ coincidence with $a\Lambda(1/48)=0.5700$ (other $\beta$ give $1.208$, $7.4\times10^3$, $5.2\times10^{10}$); the free field is equally ``constant'' ($2\pi/8=0.7854$)\\
8 & $\Phi$ (object mismatch) & an activity sum $\Phi(0)=\sum_{r\ne0}c_r\dim r=\beta=2N^2u$ & \LN & $\Phi$ is a sum over $f$ and $\bar f$ (each carrying $\dim f=N$); $\bar C$ wants one slab, and one plaquette may take $f$ or $\bar f$ but not both\\
9 & $\varphi_{\rm dec}$ & the decay density of the decoration gas, $\log Z_{\rm dec}=\varphi_{\rm dec}T+\log D_\infty+\varepsilon_T$ & \LN & Fekete already gives its rate and its existence without any positivity; a tool whose output is already a theorem adds nothing\\
\multicolumn{5}{@{}l}{\emph{Circular (9)}}\\
10 & positivity improvement & Perron--Frobenius locks the vacuum at every $\beta>0$ & \LN & assumes $r>0$: a uniform cone-improvement constant is equivalent to a uniform gap, so the hypothesis is the conclusion\\
11 & minimality & $\lambda_1(\nu)\ge\operatorname{ess\,inf}G''$ & \LN & reduces to $\operatorname{ess\,inf}G''=H1$, and the Savvidy tail \cite{savvidy,no,leutwyler} makes $\operatorname{ess\,inf}G''<0$\\
12 & operator inequalities & Temple \cite{temple}, Weyl, Lieb--Thirring, Cheeger, ELS & \LN & $w_1$ is not a function of $\mathrm{Spec}(H)$: at fixed spectrum $w_1=\cos^2\theta\in[0,1]$; the input is the target\\
13 & definability & $o$-minimality can be had on $L/a\le C$ & \LN & the non-circular version requires $(N=\infty,\beta=\infty)$ to be in the family with value $>0$, which is the Clay problem itself\\
14 & strong--weak bridge & $\beta<3.777\times10^{-3}$ gives $f_1^2\ge r_0$; $\beta\to\infty$ gives $m_{\rm lat}=D_0e^{-c\beta}$ & \LN & the two ends are of different types, $662\times$ apart; no interpolation theorem (non-negative real-analytic can be zero in between)\\
15 & $\Delta_0$ & $\exists\Delta_0>0\ \forall a<a_0:\ m_{0^{++}}(a)\ge\Delta_0$ & \LN & M6: $\Delta_{\rm dyn}>0$ while $m=0$ (not sufficient); 2D: $\Delta_{\rm dyn}=0$ while $m>0$ (not necessary); and the floor is the target renamed (\S\ref{sec:reduction-legs}(iii)), which is why the row sits in this group rather than among the counterexamples\\
16 & $T4$ & the same statement in physical units, $=H1$ dimension-corrected & \LN & its antecedent is its conclusion in another notation: $C>0\iff(T4)\iff H1$, so the step consumes what it proves\\
17 & $H1$ & a one-sided variance bound at typical field and $k\to0$ & \LN & $H1\not\Rightarrow m>0$: the thin atom witnesses a gap-free state satisfying $H1$\\
18 & $M(\beta)=m/\Lambda$ & an identity ($\gamma_m=0$): flow invariant and positive by construction & \LN & an identity cannot floor anything; its positivity is free and its uniformity is the assumption\\
\multicolumn{5}{@{}l}{\emph{Counterexamples (2)}}\\
19 & $w_1(L)$ monotonicity & the mechanism family realises both directions & \LN & $w_1=2/(L-1)$ is non-increasing with limit $0$; non-increasing does not imply $\inf>0$; identity of the mechanism family, not a defect of one member\\
20 & tight family $+$ limit atom & a compact family in $\mathcal P([0,3])$, weakly closed & \LN & $\int s\,d\mu\equiv2.4$, $E$ jumps $0.3\to0$, $\mathrm{TV}\equiv0.6$: tightness does not give a uniform floor, the atom survives only at the limit point\\
\multicolumn{5}{@{}l}{\emph{By-products (2)}}\\
21 & block-Loewner & $K_\beta\succeq K_{\beta'}$ for all $N$ & \LT & not a route to $w_1$ (blocked by the $D$-factor and by not transporting eigenvectors) but it repairs the previously unproved step of D05\\
22 & QR-sampler bug & the sampler fabricates sign changes & \LE & four corpus conventions need revision; earlier numerical ``$c_r$ changes sign'' conclusions are void\\
\end{longtable}
}

\src{Row sources, in the order of the table.
1: \code{U01.md} \S\S1--3 (\code{LOOP-RUNBOOK.md}:19064--19070).
2: \code{U03.md} (\code{LOOP-RUNBOOK.md}:19307--19311); the bound $\widehat w_1(1)=e^{1.68}\cdot0.184\approx0.83$ is \code{U14}/\code{UNIF2}.
3: \code{U11.md} \S1, \code{LOOP-RUNBOOK.md}:19257--19263.
4: \code{U12.md}; \code{LOOP-RUNBOOK.md}:19295--19298.
5: \code{U10.md} \S3 (factor $2.0975$, window $[0.690,1.448]$); acid test at $\beta=0$: \code{LOOP-RUNBOOK.md}:19301--19302.
6: \code{U13.md} \S3 (Lorentzian witness, $f''$ sign change at $2v/\sqrt3$); \code{LOOP-RUNBOOK.md}:19136--19144.
7: \code{MON3.md} \S3(iv) ($\delta_A\cdot a=0.5686159$ versus $a\Lambda(1/48)=0.5700$, ratio $0.9976$; other $\beta$: $1.208/7.394\times10^3/5.188\times10^{10}$); free field $2\pi/8=0.7854$, \code{MON3.md}:185.
8: \code{CB4.md} \S1 row 4 and \code{LOOP-RUNBOOK.md}:16589, :16622; $\Phi(0)=\beta=2N^2u$ is \code{C05-AUDITCBAR.md}:18,52.
9: \code{U10.md} \S1 (J5, threefold redundancy); Fekete side \code{RS5.md} J1, \code{P4-OTHER.md}:145.
10: \code{U10.md} \S\S2--3; \code{U15.md} row 10.
11: \code{U15.md} row 11; $\operatorname{ess\,inf}G''<0$ with $K_S(2.5)=-119$: \code{GAP2.md}:72.
12: \code{U11.md} \S2; witness \code{LOOP-RUNBOOK.md}:19241--19252.
13: \code{U04.md} \S3 (\code{LOOP-RUNBOOK.md}:19180--19194).
14: \code{U13.md} \S1--2 ($2.5/3.777\times10^{-3}=661.9$, i.e.\ $662\times$, \code{U13.md}:56); \code{LOOP-RUNBOOK.md}:19120--19133.
15: \code{LOOP-RUNBOOK.md}:18914--18916 (\S159.1); \code{Q03-TARGET.md}:127; 2D $m=3g^2/8$ and $\nu_{\frac14{\rm tr}F^2}\equiv0$: \code{N02-FOURMODEL.md}, \code{LOOP-RUNBOOK.md}:19088.
16: \code{M05-RATIO.md}:97--100; circularity \code{A04-RATIO2.md}:164 (\code{LOOP-RUNBOOK.md}:18894).
17: \code{Q6-H1-STATUS.md}:208 (thin atom); \code{R6-ATOM-ADVERSARIAL.md}:254 (gapped, atomless); \code{LOOP-RUNBOOK.md}:19048--19053.
18: \code{MON3.md} \S1; \code{LOOP-RUNBOOK.md}:19391 ($dM/d\ln a=0$, $\gamma_m=0$).
19: \code{U02.md}, \code{MON1.md} (\code{LOOP-RUNBOOK.md}:18822--18824, :19300).
20: \code{U16.md} \S3.2, table at :157--164 and :172 (\code{LOOP-RUNBOOK.md}:19351--19355).
21: \code{U09.md} \S3 (\code{LOOP-RUNBOOK.md}:19460--19464, :19490).
22: \code{U09.md} \S3 (\code{LOOP-RUNBOOK.md}:19467--19471).}

Two notes on the table, both honest bookkeeping rather than refinement of the content.
First, the three group headings are not a partition of the \emph{methods} but of the
\emph{outputs}: route 6 proves a true theorem and its hypothesis is false, and so does
route 2, which is why the direction of the bound, not the truth of the theorem, is the
discriminating property. Second, route 14 is graded \LN{} here on the strength of the
interpolation step (\code{LOOP-RUNBOOK.md}:19127--19133), while the summary ledger of that
sweep grades the leg \LC{} with the circularity as its cause
(\code{LOOP-RUNBOOK.md}:19012). The two gradings are both recorded; this paper does not
silently merge them.

\noindent\textbf{The inequalities behind the table, cited.} The correlation inequalities behind rows 2
and 10 are Griffiths \cite{griffiths}, Kelly--Sherman \cite{ksherman}, FKG \cite{fkg} and Ginibre
\cite{ginibre}; log-concavity enters through Pr\'ekopa \cite{prekopa} and Brascamp--Lieb \cite{bl};
row 12 is Temple's bound \cite{temple} together with the Weyl and Lieb--Thirring inequalities; the
``output is an upper bound'' pattern of rows 3--6 is what the transport--entropy criterion
\cite{vrs} and the convolution and entropy-power inequalities \cite{royen,gsv} also produce, and the
negative-dependence theory of \cite{bbl} is what makes a naive FKG substitute unavailable for a
lattice gauge measure; and the rates of rows 5, 7 and 9 exist as stated only inside the
strong-coupling window, which is the log-Sobolev and mixing literature of
\cite{js,sinclair,feller,bbs}. Every one of these is a true theorem. The column that matters is
\emph{what it gives}, not whether it is true.

\subsection{The three layers of the obstruction}\label{sec:survey-three}

The twenty failures are not twenty independent accidents. They decompose into three
layers, each of which is independently sufficient to stop the route, and all three point
the same way.

\paragraph{1. The instrument layer.} The weak-$*$ instruments --- moments, $\varphi^U$,
$a_n$, and every other quantity the campaign can actually compute --- cannot resolve an
atom. The decisive count is $0/14$: of fourteen candidate instruments tested against this
question, none works, and the failure is located in the class of instruments (finite data,
weak-$*$ continuous), not in infinite dimensionality. The mechanism is exact: every
computable quantity is a weak-$*$ continuous functional of $\nu$, whereas $w_1$ is the
mass of a point. A continuous functional cannot see a point mass, so no sequence of
computable quantities converges to $w_1$ in a way that could certify a floor. The one
non-weak-$*$ instrument available is $H1$, and $H1\not\Rightarrow m>0$ (row 17).

\paragraph{2. The continuity layer.} The map $\nu\mapsto\nu(\{x\})$ is only upper
semicontinuous, and a floor requires lower semicontinuity. Concretely, $w_1$ is
\emph{discontinuous} in the K\"allen--Lehmann weak topology: there is a family $\nu_\varepsilon$
with $w_1(\nu_\varepsilon)\to\varepsilon\to0$ across a jump, so that
\[
\inf\;=\;0 \qquad\text{while}\qquad w_1>0\ \text{pointwise.}
\]
Upper semicontinuity permits exactly this downward jump. The quantitative form is the same
statement: for any weak-$*$ continuous $g$ the error obeys $\sup|E-g|\ge f_1^2/2$, so a
continuous instrument cannot even approximate the atom to better than half its mass.
Consequently the target of \S\ref{sec:survey-target} \emph{may be false}, and its being
false is not a defect of the route but a property of the object. The sentence is about $\nu$ as an
unrestricted measure on $[0,\infty)$; read on the atom route of \S\ref{sec:reduction-equiv} the
same statement is the target itself, and that is the entire content of the qualifier printed there.

\paragraph{3. The implication layer.} Even granting the atom, atom-ness does not imply a
gap. The two solvable models separate the implications in opposite directions: 2D
Yang--Mills has a gap and no atom ($m=3g^2/8>0$ while $\nu_{\frac14{\rm tr}F^2}\equiv0$),
and the M6 model has an atom and no gap ($w_1=0.5$ while $m=0$). The stronger statement is
available in 2D and is the cleanest single line in the survey: there $m(L)=\sigma L\to\infty$
while $w_1=0$, so the gap \emph{diverges} while the atom fraction is identically zero.
$w_1$ and $m$ are therefore completely decoupled --- not merely non-proportional.

\begin{v3box}
\textbf{The three layers are independent, and the direction is common.} Layer 1 says the
instruments cannot see the object; layer 2 says the object's own topology forbids the
floor the target asks for; layer 3 says the floored object would not imply the goal even
if floored. No one of the three presupposes another, so removing one obstruction would
leave the other two intact. And the direction of the twenty routes is uniform:
\textbf{upper bounds are in ample supply, every lower bound is circular or
counterexampled.}
\src{\code{U15.md} (the $0/14$ count and the instrument-class localisation);
\code{U08.md} ($\nu_\varepsilon$); \code{R3} via \code{U11.md} \S3 ($\sup|E-g|\ge f_1^2/2$);
\code{U01.md} (usc versus lsc); \code{U06.md} and \code{LOOP-RUNBOOK.md}:19520--19537
(2D, $m(L)=\sigma L\to\infty$ with $w_1=0$); \code{Q03-TARGET.md}:127 and
\code{LOOP-RUNBOOK.md}:19085--19090 (M6); summary \code{LOOP-RUNBOOK.md}:19164--19170,
:19316--19325, :19569--19578.}
\end{v3box}

\subsection{What survives}\label{sec:survey-survives}

Two items survive the sweep, and one of them is a bug.

\begin{itemize}
\item \textbf{A block-Loewner inequality, \LT.} On the dual side the Griffiths/Bochner
  positivity gives $c_r(\beta)\ge0$ and $\beta$-non-decreasing (tensor-product
  multiplicities are non-negative; verified numerically on SU(3) at $\beta=2.5$ and
  $5.7$), and hence $K_\beta\succeq K_{\beta'}$ for $\beta\ge\beta'$, for all $N$. This
  repairs the previously unproved step of D05. It does not reach $w_1$: it is blocked by
  the $D$-factor and it does not transport eigenvectors.
\item \textbf{A sampler bug, \LE.} The same sweep found that the QR sampler fabricates
  sign changes. Four corpus conventions therefore require revision, and the earlier
  numerical determinations of ``$c_r$ changes sign'' are void. This is a by-product of a
  negative survey, and it is the only place in the paper where a negative result removed a
  positive claim rather than failing to establish one.
\end{itemize}

The last observation of the survey closes the FKG family as well, and closes it at $H1$.
The only monotonicity $w_1$ admits is on the three spectral points $(s_1,s_2,\bar E)$, in
the right direction:
\[
w_1\;\ge\;\frac{s_2-\bar E}{s_2-s_1},
\]
and $w_1$ is not monotone in $\beta$ (a $\tfrac14\sin^2\theta$ hump) and is constant under
$O\mapsto\lambda O$ ($\gamma=0$). The threshold form of that inequality is
\[
w_1\ge w_0 \iff \frac{\delta_1}{\varepsilon}\le 1-w_0,
\]
which is precisely the isolation of the second support point, i.e.\ $H1$. So the FKG route
does not deliver a floor either: it delivers $H1$, and $H1$ does not imply $m>0$ (row 17).
Combined with the instrument layer, the conclusion is that positivity in this family reaches
the second support point and stops there.

\src{\code{U09.md} \S3 and \code{LOOP-RUNBOOK.md}:19460--19464, :19467--19471, :19490--19498
(block-Loewner, sampler bug, and the reduction $w_1\ge w_0\iff\delta_1/\varepsilon\le1-w_0$ =
second support point $=H1$); $H1\not\Rightarrow m>0$: \code{Q6-H1-STATUS.md}:208,
\code{LOOP-RUNBOOK.md}:19048--19053.}
